\subsection{表示与格罗滕迪克群}

在本小节中，我们假设 K 是一个交换域。给定 G 的一个有限维线性表示 \((M,\pi)\)，用 \([\pi]\) 表示 K{[}G{]}-模 M 在格罗滕迪克环 \(\mathrm{R}_{\mathrm{K}}(\mathrm{G})\) 中的类（VIII, 第 198 页，例 1）。

由环 \(\mathrm{R}_{\mathrm{K}}(\mathrm{G})\) 上合成律的定义，若 \(\pi\) 与 \(\pi'\) 是 G 的有限维线性表示，则我们有

\[
[ \pi \oplus \pi^ {\prime} ] = [ \pi ] + [ \pi^ {\prime} ], \quad [ \pi \otimes \pi^ {\prime} ] = [ \pi ] [ \pi^ {\prime} ].
\]

环 \(R_{K}(G)\) 的单位元是 G 的单位表示的类（VIII, 第 399 页，例 1）。

环 \(R_{K}(G)\) 是一个自由 Z-模，以 K 上有限维单 K{[}G{]}-模的类所成的集合为基。当 \(|G|\) 在 K 中可逆时，两个有限维表示同构，当且仅当它们在 \(R_{K}(G)\) 中有相同的类（VIII, 第 190 页，推论，及第 401 页，推论 1）。

对 G 的每个有限维线性表示 \((M,\pi)\)，逆步表示 \(\pi^{\vee}\) 也是有限维的。表示 \(\pi\) 与 \((\pi^{\vee})^{\vee}\) 同构。若 \(\pi'\) 也是 G 的一个有限维线性表示，则表示 \((\pi\otimes\pi')^{\vee}\) 与 \(\pi^{\vee}\otimes\pi'^{\vee}\) 同构。若

\[
0 \longrightarrow \mathrm{M} ^ {\prime} \longrightarrow \mathrm{M} \longrightarrow \mathrm{M} ^ {\prime \prime} \longrightarrow 0
\]

是 K 上有限维 K{[}G{]}-模的一个正合列，则逆步表示给出 K{[}G{]}-模的一个正合列

\[
0 \longrightarrow \mathrm{M} ^ {\prime \prime \vee} \longrightarrow \mathrm{M} ^ {\vee} \longrightarrow \mathrm{M} ^ {\prime \vee} \longrightarrow 0
\]

（II, §7, No.~5, 第 299 页，定理 5 的推论）。由格罗滕迪克群的泛性质（VIII, 第 186 页，命题 4），存在一个自同构 \(c \mapsto c^{\vee}\)（环 \(\mathrm{R}_{\mathrm{K}}(\mathrm{G})\) 的自同构），它由 \([\pi]^{\vee} = [\pi^{\vee}]\)（对每个有限维线性表示 \(\pi\)，\(\mathrm{G}\) 的表示）刻画；我们有 \((c^{\vee})^{\vee} = c\)（对每个元素 \(c\)，\(\mathrm{R}_{\mathrm{K}}(\mathrm{G})\) 的元素）。

设 H 是 G 的一个子群。限制保持正合列。于是由格罗滕迪克群的泛性质（同前引文），存在一个群同态，记作 \(Res_{H}^{G}\)，它从 \(R_{K}(G)\) 到 \(R_{K}(H)\)，由关系式

\[
\mathrm{Res} _ {\mathrm{H}} ^ {\mathrm{G}} [ \pi ] = [ \mathrm{Res} _ {\mathrm{H}} ^ {\mathrm{G}} (\pi) ]\tag{6}
\]

（对 G 的每个有限维线性表示 \(\pi\)）刻画；它是一个环同态。若 H 在 G 中有有限指数，则由 II, §7, No.~7, 第 306 页的命题 14，我们可以类似地定义一个群同态 \(\mathrm{Ind}_{\mathrm{H}}^{\mathrm{G}}\)，从 \(\mathrm{R}_{\mathrm{K}}(\mathrm{H})\) 到 \(\mathrm{R}_{\mathrm{K}}(\mathrm{G})\)，它由关系式

\[
\mathrm{Ind} _ {\mathrm{H}} ^ {\mathrm{G}} [ \sigma ] = [ \mathrm{Ind} _ {\mathrm{H}} ^ {\mathrm{G}} (\sigma) ]\tag{7}
\]

对 H 的每个有限维线性表示 \(\sigma\) 成立。

由 VIII, 第 399 页的关系式 (1)、VIII, 第 400 页的关系式 (3) 以及上面提到的格罗滕迪克群的泛性质，存在一个环同态 \(\Theta_{\mathrm{G}}\)，从 \(\mathrm{R}_{\mathrm{K}}(\mathrm{G})\) 到中心函数的代数 \(\mathcal{Z}_{\mathrm{K}}(\mathrm{G})\)（\(\mathrm{G}\) 上中心函数的代数），它由 \(\Theta_{\mathrm{G}}([\pi]) = \chi_{\pi}\)（对每个有限维表示 \(\pi\)，\(\mathrm{G}\) 的表示）刻画。

若 H 是 G 的子群，则与群 G 和 H 相伴的同态 \(\Theta_{G}\) 与 \(\Theta_{H}\) 同运算 \(Res_{H}^{G}\) 与 \(Ind_{H}^{G}\) 相容（VIII, 第 402 页，以及 VIII, 第 404 页，命题 3）。

设 G 是两个群的乘积 \(G^{\prime} \times G^{\prime\prime}\)。由 VIII, 第 213 页注 2，存在一个 Z-线性映射 \(\kappa\)，从 \(\mathrm{R}_{\mathrm{K}}(\mathrm{G}^{\prime}) \otimes_{\mathbf{Z}} \mathrm{R}_{\mathrm{K}}(\mathrm{G}^{\prime\prime})\) 到群 \(\mathrm{R}_{\mathrm{K}}(\mathrm{G}^{\prime} \times \mathrm{G}^{\prime\prime})\)，它由关系式 \(\kappa([\pi^{\prime}] \otimes [\pi^{\prime\prime}]) = [\pi^{\prime} \boxtimes \pi^{\prime\prime}]\) 刻画，其中 \(\pi^{\prime}\)（相应地，\(\pi^{\prime\prime}\)）是 \(G^{\prime}\)（相应地，\(G^{\prime\prime}\)）的一个有限维表示，而 \(\pi^{\prime} \boxtimes \pi^{\prime\prime}\) 是外张量积（VIII, 第 400 页，例 5）。它是一个环同态。若域 K 是代数闭的，则映射 \(\kappa\) 是一个同构（VIII, 第 213 页，注 2）。

设 G 是两个群的乘积 \(G' \times G''\)。用 \(\psi\) 表示从 \(\mathcal{Z}_{\mathrm{K}}(\mathrm{G}') \otimes_{\mathrm{K}} \mathcal{Z}_{\mathrm{K}}(\mathrm{G}'' )\) 到 \(\mathcal{Z}_{\mathrm{K}}(\mathrm{G})\) 的、把 \(f' \otimes f''\) 映为函数 \((g', g'') \mapsto f'(g')f''(g'')\) 的同态。下图交换：

\[
\mathrm{R} _ {\mathrm{K}} (\mathrm{G} ^ {\prime}) \otimes_ {\mathbf {Z}} \mathrm{R} _ {\mathrm{K}} (\mathrm{G} ^ {\prime \prime}) \xrightarrow {\kappa} \mathrm{R} _ {\mathrm{K}} (\mathrm{G})
\]

\[
\bigg \downarrow \Theta_ {\mathrm{G} ^ {\prime}} \otimes \Theta_ {\mathrm{G} ^ {\prime \prime}} \qquad \qquad \bigg \downarrow \Theta_ {\mathrm{G}}
\]

\[
\mathcal {Z} _ {\mathrm{K}} (\mathrm{G} ^ {\prime}) \otimes_ {\mathrm{K}} \mathcal {Z} _ {\mathrm{K}} (\mathrm{G} ^ {\prime \prime}) \xrightarrow {\psi} \mathcal {Z} _ {\mathrm{K}} (\mathrm{G}).
\]

